\documentclass[11pt]{article}
\usepackage[utf8]{inputenc}
\usepackage[spanish,es-nodecimaldot,es-noshorthands]{babel}
\usepackage[margin=2.3cm]{geometry}
\usepackage{amsmath,amssymb}
\usepackage{enumitem}
\usepackage{tcolorbox}
\tcbuselibrary{skins,breakable}
\usepackage{pgfplots}
\pgfplotsset{compat=1.18}
\usepackage{tikz}
\usepackage{graphicx}
\usepackage{array}
\usepackage[hidelinks]{hyperref}

% Notación española opcional (descomenta si quieres \sen en lugar de \sin)
% \DeclareMathOperator{\sen}{sen}
% \DeclareMathOperator{\tg}{tg}
% \DeclareMathOperator{\ctg}{ctg}

\newtcolorbox{sol}{colback=blue!5,colframe=blue!50!black,title=Solución paso a paso,breakable}
\newtcolorbox{teo}{colback=orange!8,colframe=orange!70!black,title=Teoría,breakable}
\newtcolorbox{nota}{colback=green!5,colframe=green!50!black,title=Idea clave,breakable}
\newtcolorbox{tip}{colback=yellow!10,colframe=yellow!60!black,title=Consejo,breakable}
\newtcolorbox{resp}{colback=red!5,colframe=red!50!black,title=Respuesta,breakable}
\newtcolorbox{reg}{colback=red!5,colframe=red!60!black,title=Regla,breakable}
\newtcolorbox{paso}[1]{colback=gray!5,colframe=gray!50!black,title={#1},breakable}

\graphicspath{{../images/}{images/}}
\setlength{\parskip}{0.4em}
\setlength{\parindent}{0pt}

% ---- Formato "clase" (como Clase Leslie, 14 oct 2026) ----
\usepackage[scaled]{helvet}
\renewcommand{\familydefault}{\sfdefault}
\usepackage{sfmath}
\usepackage{bm}
% X sobre la letra de la opción correcta
\newcommand{\correcta}[1]{\makebox[0pt][l]{\hspace{-0.25em}\raisebox{-0.25ex}{\textcolor{red!75!black}{\huge$\bm{\times}$}}}#1}
% anotación a mano en color
\newcommand{\anot}[2]{\textcolor{#1}{\bfseries #2}}
% número con la operación escrita debajo (en rojo)
\newcommand{\bajo}[2]{\underset{\textcolor{red!75!black}{#2}}{#1}}
% cálculo grande en azul, centrado
\newcommand{\calculo}[1]{\begin{center}\Large\color{blue!60!black}$\displaystyle #1$\end{center}}
% figura original con anotaciones encima: coordenadas 0..1 (x desde la izquierda, y desde abajo)
\newenvironment{figanotada}[2]{\begin{center}\begin{tikzpicture}\node[anchor=south west,inner sep=0] (img) at (0,0) {\includegraphics[width=#1]{#2}};\begin{scope}[x={(img.south east)},y={(img.north west)}]}{\end{scope}\end{tikzpicture}\end{center}}
% explicación corta con la fórmula
\newtcolorbox{explica}{colback=blue!3,colframe=blue!45!black,title={¿Cómo se resuelve?},fonttitle=\bfseries,breakable}
% opciones apiladas como en el examen; la correcta se marca con \correcta{A)}
\newenvironment{opciones}{\par\smallskip\begin{list}{}{\leftmargin=2.2em\itemsep=0.35em\parsep=0pt\topsep=0pt}\item[]}{\end{list}\par}
\newcommand{\op}{\par}
% definición de teoría al estilo de la clase
\newcommand{\defi}[2]{\par\noindent\textbf{#1:} #2\par\medskip}

\usetikzlibrary{arrows.meta,positioning,calc,matrix,fit,backgrounds}
\tikzset{
  v/.style={circle,draw=black!80,thick,fill=white,minimum size=6.5mm,inner sep=0pt,font=\small},
  ar/.style={-{Stealth[length=2.3mm]},thick,blue!60!black},
  ok/.style={-{Stealth[length=2.4mm]},very thick,blue!60!black},
  bad/.style={-{Stealth[length=2.4mm]},very thick,red!75!black,dashed},
  cam/.style={-{Stealth[length=2.4mm]},very thick,orange!90!black},
  gs/.style={-{Stealth[length=2.3mm]},thick,green!45!black},
  gr/.style={-{Stealth[length=2mm]},gray!55},
  lz/.style={-{Stealth[length=2mm]},thick,blue!60!black,loop above,min distance=6mm,in=60,out=120},
  oval/.style={draw=black!70,thick,rounded corners=12pt}
}
% equivalence class with a bullet on top, as written in the exam key
\newcommand{\cl}[1]{\overset{\bullet}{#1}}
\newcommand{\R}{\mathcal R}
\newcommand{\Sc}{\mathcal S}
\newcommand{\Sb}{\overline{\mathcal S}}
% highlighted keyword
\newcommand{\hl}[1]{\colorbox{green!35}{#1}}
% justification column
\newcommand{\jus}[1]{&&\text{#1}}

\title{\textbf{Clase --- II Examen Parcial (solucionario)}\\[0.3em]\large MA-1401 Matemática Discreta --- Instituto Tecnológico de Costa Rica\\[0.3em]\normalsize I semestre 2026 --- Examen ordinario --- Puntaje total: 28 puntos}
\author{}
\date{}

\begin{document}
\maketitle
\tableofcontents

% ============================================================
\clearpage
\section*{Teoría 1: relaciones, gráfico y matriz}
\addcontentsline{toc}{section}{Teoría 1: relaciones, gráfico y matriz}

\begin{minipage}[t]{0.55\textwidth}
\defi{Relación y gráfico}{una relación $\R$ de $A$ en $B$ es una regla que dice qué elemento de $A$ ``se relaciona'' con qué elemento de $B$. Se guarda como una lista de parejas ordenadas, el \textbf{gráfico} $G_\R$:}
\[a\,\R\,b\iff (a,b)\in G_\R\subseteq A\times B.\]
\textbf{Mini‑ejemplo:} $A=\{1,2,3\}$ y $a\,\R\,b\iff a<b$. Se revisan las parejas: $1<2$, $1<3$, $2<3$ son ciertas, así que $G_\R=\{(1,2),(1,3),(2,3)\}$.
\end{minipage}\hfill
\begin{minipage}[t]{0.4\textwidth}
\vspace{0pt}\centering
\begin{tikzpicture}[scale=0.9]
  \draw[oval] (-0.6,-0.4) rectangle (0.6,2.4); \node at (0,2.75) {$A$};
  \draw[oval] (2.4,-0.4) rectangle (3.6,2.4); \node at (3,2.75) {$A$};
  \foreach \n/\y in {1/2,2/1,3/0}{\node[v] (L\n) at (0,\y) {\n}; \node[v] (D\n) at (3,\y) {\n};}
  \draw[ar] (L1) -- (D2); \draw[ar] (L1) -- (D3); \draw[ar] (L2) -- (D3);
  \node[font=\scriptsize,blue!60!black] at (1.5,-0.8) {cada flecha es un par de $G_\R$};
\end{tikzpicture}
\end{minipage}

\medskip
\begin{minipage}[t]{0.55\textwidth}
\defi{Digrafo}{cuando $\R$ va de $A$ en el mismo $A$, se dibuja cada elemento una sola vez como un punto y cada par $(a,b)$ como una flecha de $a$ hacia $b$. Un par $(a,a)$ es una flecha que sale y vuelve al mismo punto: se llama \textbf{lazo}.}
En el dibujo, la relación $G=\{(1,1),(1,2),(2,3),(3,2)\}$: el lazo en $1$ es el par $(1,1)$, y la ida y vuelta entre $2$ y $3$ son los pares $(2,3)$ y $(3,2)$.
\end{minipage}\hfill
\begin{minipage}[t]{0.4\textwidth}
\vspace{0pt}\centering
\begin{tikzpicture}
  \node[v] (a) at (0,0) {1}; \node[v] (b) at (2,0) {2}; \node[v] (c) at (1,-1.6) {3};
  \path (a) edge[lz] (a);
  \draw[ar] (a) -- (b);
  \draw[ar] (b) to[bend left=15] (c);
  \draw[ar] (c) to[bend left=15] (b);
  \node[font=\scriptsize,blue!60!black] at (-0.9,0.9) {lazo $(1,1)$};
\end{tikzpicture}
\end{minipage}

\medskip
\begin{minipage}[t]{0.55\textwidth}
\defi{Matriz asociada}{se ordena $A=\{a_1,a_2,\dots,a_n\}$ en el orden en que está escrito. Se arma una tabla de $0$ y $1$: la \textbf{fila} es el \textbf{primer} elemento del par y la \textbf{columna} es el \textbf{segundo}.}
\[M_\R[i,j]=1\iff a_i\,\R\,a_j .\]
\textbf{Gráfico $\to$ matriz:} por cada par $(a,b)$ se pone un $1$ en la fila de $a$, columna de $b$; todo lo demás es $0$.\\
\textbf{Matriz $\to$ gráfico:} se recorre fila por fila; cada $1$ da el par (nombre de la fila, nombre de la columna).
\end{minipage}\hfill
\begin{minipage}[t]{0.4\textwidth}
\vspace{0pt}\centering
\begin{tikzpicture}
  \matrix (m) [matrix of math nodes,nodes={minimum size=6mm,anchor=center},column sep=1pt,row sep=1pt]{
    & |[gray]| 1 & |[gray]| 2 & |[gray]| 3\\
    |[gray]| 1 & 0 & 1 & 1\\
    |[gray]| 2 & 0 & 0 & 1\\
    |[gray]| 3 & 0 & 0 & 0\\};
  \draw[thick] (m-2-2.north west) rectangle (m-4-4.south east);
  \draw[red!75!black,very thick] (m-2-4) circle (3.2mm);
  \draw[orange!90!black,very thick,->] (m-2-1.west) ++(-0.5,0) node[left,font=\scriptsize] {fila 1} -- (m-2-1.west);
  \draw[orange!90!black,very thick,->] (m-1-4.north) ++(0,0.5) node[above,font=\scriptsize] {columna 3} -- (m-1-4.north);
  \node[red!75!black,font=\scriptsize\bfseries,align=center] at (1.0,-1.75) {fila 1, columna 3 vale 1:\\ $1\,\R\,3$};
\end{tikzpicture}
\end{minipage}

\medskip
\textbf{¿Por qué así?} La matriz es solo otra forma de escribir la misma lista de pares: cada casilla responde la pregunta ``¿el de esta fila se relaciona con el de esta columna?'', $1$ = sí, $0$ = no. Lo único que hay que respetar es el \textbf{orden} de $A$ tal como lo escribe el enunciado, tanto en filas como en columnas.

% ============================================================
\clearpage
\section*{Teoría 2: inversa, complemento y composición}
\addcontentsline{toc}{section}{Teoría 2: inversa, complemento y composición}

\begin{minipage}[t]{0.58\textwidth}
\defi{Relación inversa}{se le da vuelta a cada par: lo que era $(a,b)$ pasa a ser $(b,a)$.}
\[G_{\R^{-1}}=\{(b,a)\mid a\,\R\,b\},\qquad M_{\R^{-1}}=(M_\R)^T.\]
\textbf{¿Por qué la transpuesta?} Darle vuelta al par cambia ``primer elemento'' por ``segundo'', o sea cambia fila por columna. Transponer es exactamente eso: la fila $i$ pasa a ser la columna $i$.\\
\textbf{Mini‑ejemplo:} si $G_\R=\{(1,2),(2,3)\}$ entonces $G_{\R^{-1}}=\{(2,1),(3,2)\}$.
\end{minipage}\hfill
\begin{minipage}[t]{0.37\textwidth}
\vspace{0pt}\centering
\begin{tikzpicture}
  \node[v] (a) at (0,0) {1}; \node[v] (b) at (1.5,0) {2}; \node[v] (c) at (3,0) {3};
  \draw[ar] (a) to[bend left=30] node[above,font=\scriptsize] {$\R$} (b);
  \draw[ar] (b) to[bend left=30] node[above,font=\scriptsize] {$\R$} (c);
  \draw[cam] (b) to[bend left=30] node[below,font=\scriptsize] {$\R^{-1}$} (a);
  \draw[cam] (c) to[bend left=30] node[below,font=\scriptsize] {$\R^{-1}$} (b);
\end{tikzpicture}
\end{minipage}

\medskip
\begin{minipage}[t]{0.58\textwidth}
\defi{Complemento de una relación}{es todo lo que \textbf{no} está en $\R$. Si $\R$ es una relación sobre $A$, el complemento son los pares de $A\times A$ que no están en $G_\R$.}
\[G_{\overline{\R}}=(A\times A)-G_\R,\qquad M_{\overline{\R}}:\ \text{cambiar cada }0\leftrightarrow1.\]
\textbf{¿Por qué?} Cada casilla dice ``sí'' ($1$) o ``no'' ($0$). El complemento contesta lo contrario en todas las casillas.
\end{minipage}\hfill
\begin{minipage}[t]{0.37\textwidth}
\vspace{0pt}\centering
\[\begin{pmatrix}1&0\\ 1&1\end{pmatrix}\ \xrightarrow{\ 0\leftrightarrow1\ }\ \begin{pmatrix}\textcolor{red!75!black}{0}&\textcolor{red!75!black}{1}\\ \textcolor{red!75!black}{0}&\textcolor{red!75!black}{0}\end{pmatrix}\]
\end{minipage}

\medskip
\begin{minipage}[t]{0.55\textwidth}
\defi{Composición}{$S\circ R$ se lee ``\textbf{primero} $R$, \textbf{luego} $S$'' (de derecha a izquierda). Hay que caminar dos pasos: de $a$ a un intermedio $b$ con $R$, y de $b$ a $c$ con $S$.}
\[a\,(S\circ R)\,c\iff(\exists b)\,[\,a\,R\,b\ \wedge\ b\,S\,c\,].\]
\textbf{En matrices:} $M_{S\circ R}=M_R\odot M_S$ (la del primer paso va a la izquierda).\\
\textbf{¿Por qué?} El par $(a,c)$ entra si existe al menos \textbf{un camino} $a\to b\to c$. No importa por cuál $b$ se pase.
\end{minipage}\hfill
\begin{minipage}[t]{0.4\textwidth}
\vspace{0pt}\centering
\begin{tikzpicture}
  \node[v] (a) at (0,0) {$a$}; \node[v] (b) at (1.8,0) {$b$}; \node[v] (c) at (3.6,0) {$c$};
  \draw[ar] (a) -- node[above,font=\scriptsize] {$R$} (b);
  \draw[gs] (b) -- node[above,font=\scriptsize] {$S$} (c);
  \draw[cam] (a) to[bend right=35] node[below,font=\scriptsize] {$S\circ R$} (c);
  \node[font=\scriptsize] at (0,0.65) {inicio}; \node[font=\scriptsize] at (1.8,0.65) {intermedio}; \node[font=\scriptsize] at (3.6,0.65) {final};
\end{tikzpicture}
\end{minipage}

\medskip
\begin{minipage}[t]{0.5\textwidth}
\defi{Producto booleano $\odot$}{para calcular la casilla (fila $i$, columna $j$) de $M_R\odot M_S$ se pone la \textbf{fila $i$} de la primera al lado de la \textbf{columna $j$} de la segunda. Si en \textbf{alguna} posición las dos tienen $1$, la casilla vale $1$; si en ninguna, vale $0$.}
\textbf{¿Por qué?} La fila $i$ de $M_R$ dice a qué intermedios llega $a_i$; la columna $j$ de $M_S$ dice qué intermedios llegan a $a_j$. Un $1$ compartido es un intermedio común: un camino.\\
\textbf{Atajo de filas:} la fila $i$ del resultado es el ``o'' ($\vee$) de las filas de $M_S$ que corresponden a los $1$ de la fila $i$ de $M_R$.
\end{minipage}\hfill
\begin{minipage}[t]{0.46\textwidth}
\vspace{0pt}\centering
\begin{tikzpicture}
  \matrix (A) [matrix of math nodes,left delimiter=(,right delimiter=),nodes={minimum size=5.5mm}]{0&1&0\\1&0&1\\0&0&0\\};
  \node (op) [right=5mm of A] {$\odot$};
  \matrix (B) [right=5mm of op,matrix of math nodes,left delimiter=(,right delimiter=),nodes={minimum size=5.5mm}]{1&0&0\\0&1&1\\1&0&0\\};
  \draw[orange!90!black,very thick,rounded corners] (A-2-1.north west) rectangle (A-2-3.south east);
  \draw[orange!90!black,very thick,rounded corners] (B-1-1.north west) rectangle (B-3-1.south east);
  \draw[red!75!black,thick] (A-2-1) circle (2.6mm); \draw[red!75!black,thick] (B-1-1) circle (2.6mm);
  \node[font=\scriptsize,align=center,red!75!black] at ($(A.south)!0.5!(B.south)+(0,-0.6)$) {fila 2 y columna 1 tienen $1$ en la\\ posición 1 $\Rightarrow$ casilla $(2,1)=1$};
\end{tikzpicture}
\end{minipage}

% ============================================================
\clearpage
\section*{Teoría 3: propiedades de una relación sobre $A$}
\addcontentsline{toc}{section}{Teoría 3: propiedades de una relación}

\begin{minipage}[t]{0.66\textwidth}
\defi{Reflexiva}{todo elemento se relaciona consigo mismo.}
\[(\forall a\in A)\ a\,\R\,a\qquad(I_n\le M_\R:\ \text{diagonal llena de }1).\]
En el digrafo: \textbf{cada} punto tiene su lazo. Basta que falte uno para que no sea reflexiva.
\end{minipage}\hfill
\begin{minipage}[t]{0.3\textwidth}
\vspace{0pt}\centering
\begin{tikzpicture}
  \node[v] (a) at (0,0) {1}; \node[v] (b) at (1.4,0) {2}; \node[v] (c) at (2.8,0) {3};
  \foreach \n in {a,b,c} \path (\n) edge[lz] (\n);
\end{tikzpicture}
\end{minipage}

\begin{minipage}[t]{0.66\textwidth}
\defi{Simétrica}{si hay flecha de ida, tiene que haber flecha de vuelta.}
\[(\forall a,b\in A)\ a\,\R\,b\Rightarrow b\,\R\,a\qquad(M_\R=M_\R^T).\]
\end{minipage}\hfill
\begin{minipage}[t]{0.3\textwidth}
\vspace{0pt}\centering
\begin{tikzpicture}
  \node[v] (a) at (0,0) {$a$}; \node[v] (b) at (2,0) {$b$};
  \draw[ar] (a) to[bend left=25] node[above,font=\scriptsize] {ida} (b);
  \draw[cam] (b) to[bend left=25] node[below,font=\scriptsize] {vuelta} (a);
\end{tikzpicture}
\end{minipage}

\begin{minipage}[t]{0.66\textwidth}
\defi{Transitiva}{si se puede ir de $a$ a $b$ y de $b$ a $c$, tiene que existir la flecha directa de $a$ a $c$ (el ``atajo'').}
\[(\forall a,b,c\in A)\ a\,\R\,b\wedge b\,\R\,c\Rightarrow a\,\R\,c\qquad(M_{\R\circ\R}\le M_\R).\]
Para mostrar que \textbf{no} es transitiva basta encontrar \textbf{un} camino $a\to b\to c$ sin el atajo $a\to c$.
\end{minipage}\hfill
\begin{minipage}[t]{0.3\textwidth}
\vspace{0pt}\centering
\begin{tikzpicture}
  \node[v] (a) at (0,0) {$a$}; \node[v] (b) at (1.2,1.3) {$b$}; \node[v] (c) at (2.4,0) {$c$};
  \draw[ar] (a) -- (b); \draw[ar] (b) -- (c);
  \draw[cam] (a) -- node[below=2pt,font=\scriptsize] {atajo} (c);
\end{tikzpicture}
\end{minipage}

\begin{minipage}[t]{0.66\textwidth}
\defi{Antisimétrica}{entre dos elementos \textbf{distintos} nunca hay ida y vuelta (los lazos sí se permiten).}
\[(\forall a,b\in A)\ a\,\R\,b\wedge b\,\R\,a\Rightarrow a=b\qquad(M_\R\wedge M_\R^T\le I_n).\]
\end{minipage}\hfill
\begin{minipage}[t]{0.3\textwidth}
\vspace{0pt}\centering
\begin{tikzpicture}
  \node[v] (a) at (0,0) {$a$}; \node[v] (b) at (2,0) {$b$};
  \draw[ar] (a) to[bend left=25] (b);
  \draw[bad] (b) to[bend left=25] node[below,font=\scriptsize] {prohibida} (a);
\end{tikzpicture}
\end{minipage}

\begin{minipage}[t]{0.66\textwidth}
\defi{Total}{cualquier par de elementos está conectado al menos en un sentido.}
\[(\forall a,b\in A)\ a\,\R\,b\vee b\,\R\,a\qquad(M_\R\vee M_\R^T=\text{todo }1).\]
Para mostrar que \textbf{no} es total basta encontrar dos elementos sin flecha en ningún sentido.
\end{minipage}\hfill
\begin{minipage}[t]{0.3\textwidth}
\vspace{0pt}\centering
\begin{tikzpicture}
  \node[v] (a) at (0,0) {$a$}; \node[v] (b) at (2,0) {$b$}; \node[v] (c) at (1,-1.2) {$c$};
  \foreach \n in {a,b} \path (\n) edge[lz] (\n);
  \path (c) edge[-{Stealth[length=2mm]},thick,blue!60!black,loop right,min distance=6mm] (c);
  \draw[ar] (a) -- (b); \draw[ar] (c) -- (a); \draw[ar] (b) -- (c);
  \node[font=\scriptsize] at (1,-1.85) {todos conectados};
\end{tikzpicture}
\end{minipage}

\bigskip
\defi{Cómo se demuestra un ``para todo''}{se escribe primero \textbf{HQD} (hay que demostrar) con la definición de la propiedad. Luego se toman los elementos ``\textbf{fijos pero arbitrarios}'': son elementos concretos, pero sin escoger ninguno en especial, así que lo que se pruebe para ellos sirve para todos. Se termina con una cadena de $\Rightarrow$, cada una con su justificación a la derecha (Def. de $\R$, Sust., Op., Factor, \dots).}
\defi{Cómo se demuestra que una propiedad NO se cumple}{con un \textbf{contraejemplo}: números concretos que cumplen la parte ``si'' de la definición pero no la parte ``entonces''. Se verifican los dos lados con cuentas.}

% ============================================================
\clearpage
\section*{Teoría 4: relación de equivalencia y clases}
\addcontentsline{toc}{section}{Teoría 4: relación de equivalencia y clases}

\begin{minipage}[t]{0.55\textwidth}
\defi{Relación de equivalencia}{es una relación que es reflexiva, simétrica y transitiva a la vez. Agrupa los elementos que ``se parecen'' según la regla.}
\defi{Clase de equivalencia}{la clase de $a$ es el grupo de todos los que se relacionan con $a$:}
\[\cl{a}=\{x\in A : x\,\R\,a\}.\]
\defi{Conjunto cociente}{es el conjunto de todas las clases distintas:}
\[A/\R=\{\cl{a} : a\in A\}.\]
\textbf{¿Por qué forman ``cajas'' separadas?} Por la simetría y la transitividad, si dos elementos están en la misma clase, todos los de esa clase se relacionan entre sí; y dos clases distintas no comparten nada. Así, las clases parten $A$ en cajas sin elementos repetidos ni sobrantes.
\end{minipage}\hfill
\begin{minipage}[t]{0.4\textwidth}
\vspace{0pt}\centering
\begin{tikzpicture}
  \draw[oval] (-0.3,-0.3) rectangle (4.3,2.3); \node at (2,2.6) {$A$};
  \draw[blue!60!black,thick,fill=blue!8,rounded corners] (0,0) rectangle (1.2,2);
  \draw[orange!90!black,thick,fill=orange!10,rounded corners] (1.4,0) rectangle (2.6,2);
  \draw[green!45!black,thick,fill=green!10,rounded corners] (2.8,0) rectangle (4,2);
  \node at (0.6,1.4) {$a$}; \node at (0.6,0.6) {$a'$};
  \node at (2,1.4) {$b$}; \node at (2,0.6) {$b'$};
  \node at (3.4,1.0) {$c$};
  \node[font=\scriptsize] at (0.6,-0.6) {$\cl a$}; \node[font=\scriptsize] at (2,-0.6) {$\cl b$}; \node[font=\scriptsize] at (3.4,-0.6) {$\cl c$};
  \node[font=\scriptsize,align=center] at (2,-1.3) {$A/\R=\{\cl a,\cl b,\cl c\}$};
\end{tikzpicture}
\end{minipage}

\medskip
\textbf{Mini‑ejemplo:} en $A=\{1,2,3,4\}$ con $a\,\R\,b\iff a-b$ es par, la clase de $1$ son los $x$ con $x-1$ par: $\cl1=\{1,3\}$; la de $2$: $\cl2=\{2,4\}$. Como $\cl3=\cl1$ y $\cl4=\cl2$, el cociente tiene dos cajas: $A/\R=\{\{1,3\},\{2,4\}\}$.

% ============================================================
\clearpage
\section*{Pregunta 1 (forma 1: mediante gráficos)}
\addcontentsline{toc}{section}{Pregunta 1}

\textbf{\# 1. [4 puntos]} Sea $A=\{2,3,6\}$. Sean $\R$ y $\Sc$ dos relaciones definidas sobre $A$ por:
\[a\,\R\,b\iff[b=3\ \vee\ a+b=5],\qquad M_\Sc=\begin{pmatrix}1&0&1\\0&1&0\\1&1&0\end{pmatrix}.\]
Determine el \hl{gráfico y la matriz asociada} de la relación
\[\colorbox{green!35}{$\Sb\circ\R^{-1}$}.\]

\begin{center}
\begin{tikzpicture}[yscale=0.95]
  \foreach \c/\x/\t in {X/0/x,Y/3.2/y,Z/6.4/z}{
    \draw[oval] (\x-0.55,-0.5) rectangle (\x+0.55,2.5); \node at (\x,2.85) {$A$ (\t)};
    \foreach \n/\y in {2/2,3/1,6/0} \node[v] (\c\n) at (\x,\y) {\n};}
  % R^{-1}: (2,3),(3,2),(3,3),(3,6)
  \draw[ar] (X2) -- (Y3); \draw[ar] (X3) -- (Y2); \draw[ar] (X3) -- (Y3); \draw[ar] (X3) -- (Y6);
  % complement of S: (2,3),(3,2),(3,6),(6,6)
  \draw[gs] (Y2) -- (Z3); \draw[gs] (Y3) -- (Z2); \draw[gs] (Y3) -- (Z6); \draw[gs] (Y6) -- (Z6);
  \node[blue!60!black,font=\small\bfseries] at (1.6,-0.85) {$\R^{-1}$ (primero)};
  \node[green!45!black,font=\small\bfseries] at (4.8,-0.85) {$\Sb$ (luego)};
  \node[red!75!black,font=\scriptsize] at (-1.25,0) {sin flechas};
\end{tikzpicture}
\end{center}

\calculo{G_{\Sb\circ\R^{-1}}=\{\bajo{(2,2)}{2\to3\to2},\ \bajo{(2,6)}{2\to3\to6},\ \bajo{(3,2)}{3\to3\to2},\ \bajo{(3,3)}{3\to2\to3},\ \bajo{(3,6)}{3\to3\to6}\}}

\begin{explica}
\textbf{Fórmula:} composición (Teoría 2): $x\,(\Sb\circ\R^{-1})\,z\iff(\exists y\in A)\,[\,x\,\R^{-1}\,y\ \wedge\ y\,\Sb\,z\,]$. Primero se da un paso con $\R^{-1}$ y luego uno con $\Sb$ (la barra es el \textbf{complemento} de $\Sc$).

\textbf{Paso 1, gráfico de $\R$:} se revisan los $9$ pares de $A\times A$. Con $b=3$ entran $(2,3)$, $(3,3)$, $(6,3)$. Con $a+b=5$ entran $(2,3)$ ($2+3=5$) y $(3,2)$ ($3+2=5$); ninguna otra suma da $5$ ($2+2=4$, $2+6=8$, $3+6=9$, $6+6=12$, \dots). Así
\[G_\R=\{(3,2),(2,3),(3,3),(6,3)\}.\]

\textbf{Paso 2, inversa:} por Def. de inversa se le da vuelta a cada par:
\[G_{\R^{-1}}=\{(2,3),(3,2),(3,3),(3,6)\}\quad(\text{flechas azules}).\]

\textbf{Paso 3, gráfico de $\Sc$:} se lee la matriz fila por fila con el orden $2,3,6$. Fila de $2$: $(1,0,1)$, unos en las columnas $2$ y $6$, da $(2,2),(2,6)$. Fila de $3$: $(0,1,0)$, da $(3,3)$. Fila de $6$: $(1,1,0)$, da $(6,2),(6,3)$.
\[G_\Sc=\{(2,2),(2,6),(3,3),(6,2),(6,3)\}.\]

\textbf{Paso 4, complemento:} $A\times A$ tiene $3\cdot3=9$ pares; se quitan los $5$ de $G_\Sc$ y quedan $9-5=4$:
\[G_{\Sb}=\{(2,3),(3,2),(3,6),(6,6)\}\quad(\text{flechas verdes}).\]

\textbf{Paso 5, caminos de dos pasos:} se sigue cada flecha azul y después cada flecha verde que sale de donde se llegó.
\begin{center}
\renewcommand{\arraystretch}{1.25}
\begin{tabular}{|c|c|c|c|}\hline
$x$ & $y$ con $x\,\R^{-1}\,y$ & $z$ con $y\,\Sb\,z$ & pares $(x,z)$\\\hline
$2$ & $3$ & $2,\ 6$ & $(2,2),(2,6)$\\
$3$ & $2$ & $3$ & $(3,3)$\\
$3$ & $3$ & $2,\ 6$ & $(3,2),(3,6)$\\
$3$ & $6$ & $6$ & $(3,6)$ (repetido)\\
$6$ & ninguno & --- & ninguno\\\hline
\end{tabular}
\end{center}
Del $6$ no sale ninguna flecha azul porque ningún par de $G_{\R^{-1}}$ empieza con $6$; por eso la fila de $6$ queda vacía.

\textbf{Paso 6, matriz:} con el orden $2,3,6$ se pone un $1$ por cada par:
\[M_{\Sb\circ\R^{-1}}=\begin{pmatrix}1&0&1\\1&1&1\\0&0&0\end{pmatrix}.\]
\end{explica}

% ============================================================
\clearpage
\section*{Pregunta 1 (forma 2: mediante matrices)}

\calculo{M_{\Sb\circ\R^{-1}}=\bajo{\begin{pmatrix}0&1&0\\1&1&1\\0&0&0\end{pmatrix}}{(M_\R)^T}\odot\bajo{\begin{pmatrix}0&1&0\\1&0&1\\0&0&1\end{pmatrix}}{M_{\Sb}:\ 0\leftrightarrow1}=\begin{pmatrix}1&0&1\\1&1&1\\0&0&0\end{pmatrix}}

\begin{explica}
\textbf{Fórmula:} $M_{\Sb\circ\R^{-1}}=M_{\R^{-1}}\odot M_{\Sb}=(M_\R)^T\odot M_{\Sb}$. La matriz del \textbf{primer} paso ($\R^{-1}$) va a la izquierda (Teoría 2).

\textbf{Paso 1, $M_\R$:} con el orden $2,3,6$ y $G_\R=\{(3,2),(2,3),(3,3),(6,3)\}$: la fila de $2$ tiene un $1$ en la columna de $3$; la fila de $3$ tiene $1$ en las columnas de $2$ y de $3$; la fila de $6$ tiene un $1$ en la columna de $3$.
\[M_\R=\begin{pmatrix}0&1&0\\1&1&0\\0&1&0\end{pmatrix}.\]

\textbf{Paso 2, transpuesta:} la columna $j$ de $M_\R$ pasa a ser la fila $j$. Columna de $2$: $(0,1,0)$; columna de $3$: $(1,1,1)$; columna de $6$: $(0,0,0)$. Esas son las filas de $M_{\R^{-1}}$.

\textbf{Paso 3, $M_{\Sb}$:} se cambia cada $0$ por $1$ y cada $1$ por $0$ en $M_\Sc$: $(1,0,1)\to(0,1,0)$, $(0,1,0)\to(1,0,1)$, $(1,1,0)\to(0,0,1)$.

\textbf{Paso 4, producto booleano por filas:}
\begin{itemize}[leftmargin=1.3em,itemsep=2pt]
\item Fila de $2$: $(0,1,0)$ tiene un solo $1$, en la posición de $3$. Se copia la fila de $3$ de $M_{\Sb}$: $(1,0,1)$.
\item Fila de $3$: $(1,1,1)$ tiene $1$ en las tres posiciones. Se hace el ``o'' de las tres filas de $M_{\Sb}$: $(0,1,0)\vee(1,0,1)\vee(0,0,1)=(1,1,1)$.
\item Fila de $6$: $(0,0,0)$ no tiene ningún $1$, así que queda $(0,0,0)$.
\end{itemize}

\textbf{Paso 5, gráfico:} cada $1$ del resultado es un par (fila, columna):
\[G_{\Sb\circ\R^{-1}}=\{(2,2),(2,6),(3,2),(3,3),(3,6)\}.\]

\textbf{¿Tiene sentido?} Las dos formas dan lo mismo. Comprobación del par $(2,2)$ con la definición: hace falta un $y$ con $2\,\R^{-1}\,y$ y $y\,\Sb\,2$. Con $y=3$: $2\,\R^{-1}\,3$ porque $3\,\R\,2$ ($3+2=5$), y $3\,\Sb\,2$ porque en $M_\Sc$ la fila de $3$, columna de $2$ vale $0$, o sea $(3,2)\notin G_\Sc$. Sí entra.
\end{explica}

\begin{nota}
La barra sobre $\Sc$ es el complemento y cambia todo el resultado. Si se olvida y se calcula $\Sc\circ\R^{-1}$, se obtiene $(M_\R)^T\odot M_\Sc$ con filas $(0,1,0),(1,1,1),(0,0,0)$ y gráfico $\{(2,3),(3,2),(3,3),(3,6)\}$, que es \textbf{otra} relación. La clave oficial usa $M_{\Sb}$ y coincide con esta solución.
\end{nota}

\textbf{Respuesta:} $G_{\Sb\circ\R^{-1}}=\{(2,2),(2,6),(3,2),(3,3),(3,6)\}$ y $M_{\Sb\circ\R^{-1}}=\begin{pmatrix}1&0&1\\1&1&1\\0&0&0\end{pmatrix}$.

% ============================================================
\clearpage
\section*{Pregunta 2}
\addcontentsline{toc}{section}{Pregunta 2}

\textbf{\# 2. [3 puntos]} Sea $A=\{1,2,3,4,5\}$. Se define la relación $\R$ sobre $A$ mediante:
\[G_\R=\{(1,1),(2,2),(3,3),(4,4),(5,5),(1,2),(1,4),(1,3),(2,3),(2,5),(4,5)\}.\]
Para la relación $\R$, analice si satisface o no cada una de las siguientes propiedades: \hl{reflexiva}, \hl{transitiva} y \hl{total}. Justifique claramente cada respuesta.

\begin{center}
\begin{tikzpicture}[scale=1.15]
  \node[v] (1) at (90:2) {1}; \node[v] (2) at (18:2) {2}; \node[v] (3) at (-54:2) {3};
  \node[v] (4) at (-126:2) {4}; \node[v] (5) at (162:2) {5};
  \path (1) edge[lz] (1);
  \path (2) edge[-{Stealth[length=2mm]},thick,blue!60!black,loop right,min distance=6mm] (2);
  \path (3) edge[-{Stealth[length=2mm]},thick,blue!60!black,loop below,min distance=6mm] (3);
  \path (4) edge[-{Stealth[length=2mm]},thick,blue!60!black,loop below,min distance=6mm] (4);
  \path (5) edge[-{Stealth[length=2mm]},thick,blue!60!black,loop left,min distance=6mm] (5);
  \draw[ar] (1) -- (2); \draw[ar] (1) -- (3); \draw[ar] (2) -- (3);
  \draw[cam] (1) -- (4); \draw[cam] (4) -- (5); \draw[ar] (2) -- (5);
  \draw[bad] (1) to[bend right=25] node[above left,font=\scriptsize\bfseries] {falta $(1,5)$} (5);
\end{tikzpicture}
\end{center}

\calculo{\bajo{(1,4)}{\in G_\R}\ \wedge\ \bajo{(4,5)}{\in G_\R}\ \ \text{pero}\ \ \bajo{(1,5)}{\notin G_\R}}

\begin{explica}
\textbf{Paso 1, el digrafo:} se dibuja un punto por cada elemento de $A$ y una flecha por cada par de $G_\R$ (Teoría 1). Los cinco pares $(a,a)$ son lazos.

\textbf{Paso 2, reflexiva:} la definición pide $a\,\R\,a$ para \textbf{todo} $a\in A$. En $G_\R$ están $(1,1),(2,2),(3,3),(4,4),(5,5)$, uno por cada elemento de $A$: cada punto tiene su lazo.

\textbf{Paso 3, transitiva:} se busca un camino de dos flechas $a\to b\to c$ sin el atajo $a\to c$. El camino naranja $1\to4\to5$ usa $(1,4)$ y $(4,5)$, que están en $G_\R$, pero el atajo $(1,5)$ no está (flecha roja punteada). Basta un caso para que falle. (También falla con $1\to2\to5$.)

\textbf{Paso 4, total:} la definición pide que para \textbf{cualquier} par $a,b$ haya flecha en algún sentido. Entre $1$ y $5$ no hay ninguna: ni $(1,5)$ ni $(5,1)$ están en $G_\R$.

\textbf{¿Tiene sentido?} Todas las flechas entre puntos distintos van ``hacia adelante'' (de un número menor a uno mayor), nunca de regreso; por eso la relación es antisimétrica, como dice la clave, aunque no se pregunta.
\end{explica}

\textbf{Reflexiva:} Sí es reflexiva, ya que en $G_\R$ aparecen todos los pares de la forma $(a,a)$:
\[(1,1),(2,2),(3,3),(4,4),(5,5)\in G_\R.\]
\textbf{Transitiva:} No es transitiva, pues $(1,4)\in G_\R$ y $(4,5)\in G_\R$, pero $(1,5)\notin G_\R$.\\
\textbf{Total:} No es total pues $(1,5)\notin G_\R\wedge(5,1)\notin G_\R$.

\textbf{Respuesta:} $\R$ es reflexiva; no es transitiva; no es total.

% ============================================================
\clearpage
\section*{Pregunta 3a}
\addcontentsline{toc}{section}{Pregunta 3a}

\textbf{\# 3.} Sea $\R$ una relación sobre $\mathbb Z$ definida por:
\[a\,\R\,b\iff\exists k\in\mathbb Z^+\ \text{tal que}\ a^2-b^2=2k.\]
a) \textbf{[3 puntos]} Demuestre que la relación $\R$ es \hl{transitiva}.

\begin{center}
\begin{tikzpicture}
  \node[v,minimum size=8mm] (a) at (0,0) {$5$}; \node[v,minimum size=8mm] (b) at (2.4,1.6) {$3$}; \node[v,minimum size=8mm] (c) at (4.8,0) {$1$};
  \draw[ar] (a) -- node[above left,font=\scriptsize] {$25-9=2\cdot\mathbf 8$} (b);
  \draw[ar] (b) -- node[above right,font=\scriptsize] {$9-1=2\cdot\mathbf 4$} (c);
  \draw[cam] (a) -- node[below,font=\scriptsize] {$25-1=2\cdot\mathbf{12}=2(8+4)$} (c);
\end{tikzpicture}
\end{center}

\calculo{a^2-c^2=\bajo{(a^2-b^2)}{2k_1}+\bajo{(b^2-c^2)}{2k_2}=2(\bajo{k_1+k_2}{k_3\in\mathbb Z^+})}

\begin{explica}
\textbf{Definición que se usa:} transitiva (Teoría 3): $a\,\R\,b\wedge b\,\R\,c\Rightarrow a\,\R\,c$. Aquí ``$a\,\R\,b$'' significa: $a^2-b^2$ es igual a $2$ por un entero \textbf{positivo}.

\textbf{Paso 1, lo que se tiene:} $a^2-b^2=2k_1$ y $b^2-c^2=2k_2$, con $k_1,k_2$ enteros positivos.

\textbf{Paso 2, lo que se busca:} escribir $a^2-c^2$ como $2$ por un entero positivo.

\textbf{Paso 3, la idea:} al sumar las dos igualdades, el $b^2$ aparece una vez sumando y otra restando, así que se cancela: $(a^2-b^2)+(b^2-c^2)=a^2-c^2$. Del otro lado queda $2k_1+2k_2=2(k_1+k_2)$.

\textbf{Paso 4, el nuevo $k$:} $k_3=k_1+k_2$ es entero positivo porque la suma de dos positivos es positiva.

\textbf{¿Tiene sentido?} En el dibujo: $5\,\R\,3$ con $k_1=8$ y $3\,\R\,1$ con $k_2=4$; directo, $25-1=24=2\cdot12$ y $12=8+4$.
\end{explica}

\textbf{HQD:} $\R$ es transitiva $\equiv(\forall a,b,c\in\mathbb Z)\big[a\R b\wedge b\R c\Rightarrow a\R c\big]$.

Sean $a,b,c\in\mathbb Z$ fijos pero arbitrarios tales que $a\R b\wedge b\R c$.
\begin{align*}
a\R b\wedge b\R c&\Rightarrow(\exists k_1\in\mathbb Z^+)\big[a^2-b^2=2k_1\big]\wedge(\exists k_2\in\mathbb Z^+)\big[b^2-c^2=2k_2\big]\jus{Def. de $\R$}\\
&\Rightarrow(\exists k_1,k_2\in\mathbb Z^+)\big[a^2-b^2=2k_1\wedge b^2-c^2=2k_2\big]\\
&\Rightarrow(\exists k_1,k_2\in\mathbb Z^+)\big[(a^2-b^2)+(b^2-c^2)=2k_1+2k_2\big]\jus{Suma término a término}\\
&\Rightarrow(\exists k_1,k_2\in\mathbb Z^+)\big[a^2-c^2=2(k_1+k_2)\big]\jus{Simplificación}\\
&\Rightarrow(\exists k_3\in\mathbb Z^+)\big[a^2-c^2=2k_3\big],\quad k_3=k_1+k_2\jus{Def. $k_3$}\\
&\Rightarrow a\R c\jus{Def. de $\R$}
\end{align*}
Por lo tanto, $\R$ es transitiva.

\textbf{Respuesta:} $\R$ es transitiva, con $k_3=k_1+k_2\in\mathbb Z^+$.

% ============================================================
\clearpage
\section*{Pregunta 3b}
\addcontentsline{toc}{section}{Pregunta 3b}

b) \textbf{[1 punto]} Mediante un \hl{contraejemplo}, demuestre que $\R$ \hl{no es simétrica}.

\begin{center}
\begin{tikzpicture}
  \node[v,minimum size=8mm] (a) at (0,0) {$3$}; \node[v,minimum size=8mm] (b) at (4,0) {$1$};
  \draw[ok] (a) to[bend left=20] node[above,font=\scriptsize] {$9-1=8=2\cdot4$, \ $4\in\mathbb Z^+$} (b);
  \draw[bad] (b) to[bend left=20] node[below=3pt,font=\scriptsize] {$1-9=-8=2\cdot(-4)$, \ $-4\notin\mathbb Z^+$} (a);
  \node[red!75!black,font=\large\bfseries] at (2,-0.36) {$\times$};
\end{tikzpicture}
\end{center}

\calculo{3^2-1^2=\bajo{8}{=2\cdot 4}\quad\text{pero}\quad 1^2-3^2=\bajo{-8}{=2\cdot(-4),\ -4\notin\mathbb Z^+}}

\begin{explica}
\textbf{Definición que se usa:} simétrica (Teoría 3): $a\R b\Rightarrow b\R a$. Para que \textbf{no} lo sea basta un par con flecha de ida y sin flecha de vuelta.

\textbf{Paso 1, la idea:} $a^2-b^2$ y $b^2-a^2$ son el mismo número con signo cambiado. Si uno es positivo, el otro es negativo, y un negativo nunca es $2$ por un entero positivo.

\textbf{Paso 2, ida:} con $a=3$, $b=1$: $3^2-1^2=9-1=8=2\cdot4$, y $4\in\mathbb Z^+$. Así $3\R1$.

\textbf{Paso 3, vuelta:} $1^2-3^2=1-9=-8$. El único $k$ con $2k=-8$ es $k=-4$, que no es positivo. Así $1\not\R 3$.

\textbf{¿Tiene sentido?} Sí: la ida cumple la definición y la vuelta no, que es justo lo que pide un contraejemplo de simetría.
\end{explica}

Se mostrará un contraejemplo. Sean $a=3$ y $b=1$. Entonces:
\[a^2-b^2=9-1=8=2\cdot4,\]
por lo que $3\R1$. Ahora:
\[b^2-a^2=1-9=-8,\]
y no existe $k\in\mathbb Z^+$ tal que $-8=2k$. Por lo tanto, $1\not\R3$. Así, $\R$ no es simétrica.

\textbf{Respuesta:} $3\R1$ pero $1\not\R3$; $\R$ no es simétrica.

% ============================================================
\clearpage
\section*{Pregunta 4a}
\addcontentsline{toc}{section}{Pregunta 4a}

\textbf{\# 4.} Sea $A=\{0,1,2,3,4,5,6,7\}$. Se define la relación $\R$ sobre $A$ por:
\[a\,\R\,b\iff\exists k\in\mathbb Z\ t.q.\ a-b=3k.\]
$\R$ es reflexiva y transitiva (no se debe probar esto).

a) \textbf{[2 puntos]} Justifique que $\R$ es \hl{simétrica}.

\begin{center}
\begin{tikzpicture}[scale=0.82]
  \newcommand{\tri}[4]{%
    \node[v] (#1a) at ($(#4)+(0,1.5)$) {#1}; \node[v] (#1b) at ($(#4)+(-1.5,0)$) {#2}; \node[v] (#1c) at ($(#4)+(1.5,0)$) {#3};
    \path (#1a) edge[lz] (#1a);
    \path (#1b) edge[-{Stealth[length=2mm]},thick,blue!60!black,loop left,min distance=6mm] (#1b);
    \path (#1c) edge[-{Stealth[length=2mm]},thick,blue!60!black,loop right,min distance=6mm] (#1c);
    \draw[ar] (#1a) to[bend right=12] (#1b); \draw[cam] (#1b) to[bend right=12] (#1a);
    \draw[ar] (#1a) to[bend left=12] (#1c); \draw[cam] (#1c) to[bend left=12] (#1a);
    \draw[ar] (#1b) to[bend left=12] (#1c); \draw[cam] (#1c) to[bend left=12] (#1b);}
  \tri{3}{0}{6}{0,0}
  \tri{4}{1}{7}{5.8,0}
  \node[v] (2) at (10.4,0.4) {2}; \node[v] (5) at (12.2,0.4) {5};
  \path (2) edge[lz] (2); \path (5) edge[lz] (5);
  \draw[ar] (2) to[bend left=15] (5); \draw[cam] (5) to[bend left=15] (2);
\end{tikzpicture}
\end{center}

\calculo{b-a=-(a-b)=-3k=3(\bajo{-k}{k_1\in\mathbb Z})}

\begin{explica}
\textbf{Definición que se usa:} simétrica (Teoría 3): $a\R b\Rightarrow b\R a$. Aquí ``$a\R b$'' significa que $a-b$ es $3$ por un entero (un múltiplo de $3$).

\textbf{Paso 1, lo que se tiene:} $a-b=3k$ para algún entero $k$.

\textbf{Paso 2, lo que se busca:} escribir $b-a$ como $3$ por un entero.

\textbf{Paso 3, la idea:} $b-a$ es la misma resta al revés, o sea $b-a=-(a-b)$. Si $a-b=3k$, entonces $b-a=-3k=3(-k)$.

\textbf{Paso 4, el nuevo $k$:} $k_1=-k$ es entero porque el opuesto de un entero es entero. (Aquí $k$ puede ser negativo, a diferencia de la Pregunta 3, porque la definición dice $k\in\mathbb Z$.)

\textbf{¿Tiene sentido?} En el digrafo, cada flecha azul entre puntos distintos tiene su flecha naranja de regreso: por ejemplo $6-0=6=3\cdot2$ y $0-6=-6=3\cdot(-2)$.
\end{explica}

\textbf{HQD:} $\R$ es simétrica $\equiv(\forall a,b\in A)\big[a\R b\Rightarrow b\R a\big]$.

Sean $a,b\in A$ fijos pero arbitrarios tales que $a\R b$.
\begin{align*}
a\R b&\Rightarrow(\exists k\in\mathbb Z)[a-b=3k]\jus{Def. de $\R$}\\
&\Rightarrow b-a=-(a-b)\jus{Op.}\\
&\Rightarrow b-a=-3k\jus{Sust.}\\
&\Rightarrow b-a=3(-k)\jus{Factor}\\
&\Rightarrow(\exists k_1\in\mathbb Z)[b-a=3k_1],\ k_1=-k\jus{$-k\in\mathbb Z$}\\
&\Rightarrow b\R a\jus{Def. de $\R$}
\end{align*}
Por lo tanto, $\R$ es simétrica.

\textbf{Respuesta:} $\R$ es simétrica, con $k_1=-k\in\mathbb Z$.

% ============================================================
\clearpage
\section*{Pregunta 4b}
\addcontentsline{toc}{section}{Pregunta 4b}

b) \textbf{[2 puntos]} Calcule todas las \hl{clases de equivalencia} y determine, \hl{por extensión}, el conjunto cociente $A/\R$.

\begin{center}
\begin{tikzpicture}
  \draw[oval] (-0.4,-0.4) rectangle (8.4,1.4); \node at (4,1.75) {$A$};
  \draw[blue!60!black,thick,fill=blue!8,rounded corners] (0,0) rectangle (2.8,1);
  \draw[orange!90!black,thick,fill=orange!10,rounded corners] (3.0,0) rectangle (5.8,1);
  \draw[green!45!black,thick,fill=green!10,rounded corners] (6.0,0) rectangle (8.0,1);
  \foreach \n/\x in {0/0.5,3/1.4,6/2.3,1/3.5,4/4.4,7/5.3,2/6.6,5/7.4} \node at (\x,0.5) {$\n$};
  \node[font=\small] at (1.4,-0.75) {$\cl0$: residuo $0$}; \node[font=\small] at (4.4,-0.75) {$\cl1$: residuo $1$}; \node[font=\small] at (7,-0.75) {$\cl2$: residuo $2$};
\end{tikzpicture}
\end{center}

\calculo{A/\R=\big\{\bajo{\{0,3,6\}}{\cl0},\ \bajo{\{1,4,7\}}{\cl1},\ \bajo{\{2,5\}}{\cl2}\big\}}

\begin{explica}
\textbf{Definición que se usa:} clase de equivalencia (Teoría 4): $\cl a=\{x\in A: x\R a\}$, o sea todos los $x$ de $A$ tales que $x-a$ es múltiplo de $3$. Como $\R$ es reflexiva, transitiva (dato) y simétrica (parte a), es de equivalencia.

\textbf{Paso 1, tabla de restas:} para cada $x\in A$ se calcula $x-0$, $x-1$ y $x-2$ y se marca cuándo es múltiplo de $3$ (es decir, $3k$ con $k\in\mathbb Z$).
\begin{center}
\renewcommand{\arraystretch}{1.2}
\begin{tabular}{|c|cccccccc|}\hline
$x$ & 0&1&2&3&4&5&6&7\\\hline
$x-0$ & \textbf{0}&1&2&\textbf{3}&4&5&\textbf{6}&7\\
$x-1$ & $-1$&\textbf{0}&1&2&\textbf{3}&4&5&\textbf{6}\\
$x-2$ & $-2$&$-1$&\textbf{0}&1&2&\textbf{3}&4&5\\\hline
\end{tabular}
\end{center}
En negrita los múltiplos de $3$: $0=3\cdot0$, $3=3\cdot1$, $6=3\cdot2$.

\textbf{Paso 2, las clases:} de la fila $x-0$ salen $0,3,6$; de la fila $x-1$ salen $1,4,7$; de la fila $x-2$ salen $2,5$.

\textbf{Paso 3, ¿hay más clases?} No: ya aparecieron los $8$ elementos de $A$ ($3+3+2=8$), cada uno en una sola caja. Por ejemplo $\cl3=\cl6=\cl0$, $\cl4=\cl7=\cl1$ y $\cl5=\cl2$, son las mismas cajas con otro nombre.

\textbf{¿Tiene sentido?} Coincide con los tres bloques del digrafo de la parte a): dentro de cada bloque todos se relacionan y entre bloques no hay flechas.
\end{explica}

\[\cl0=\{x\in A: x-0\ \text{es múltiplo de }3\}=\{0,3,6\}.\]
\[\cl1=\{x\in A: x-1\ \text{es múltiplo de }3\}=\{1,4,7\}.\]
\[\cl2=\{x\in A: x-2\ \text{es múltiplo de }3\}=\{2,5\}.\]
Por lo tanto: $A/\R=\big\{\{0,3,6\},\{1,4,7\},\{2,5\}\big\}$.

\textbf{Respuesta:} $\cl0=\{0,3,6\}$, $\cl1=\{1,4,7\}$, $\cl2=\{2,5\}$ y $A/\R=\big\{\{0,3,6\},\{1,4,7\},\{2,5\}\big\}$.

% ============================================================
\clearpage
\section*{Teoría 5: funciones, imagen e imagen inversa}
\addcontentsline{toc}{section}{Teoría 5: funciones}

\begin{minipage}[t]{0.55\textwidth}
\defi{Función}{$f:A\to B$ es una relación en la que \textbf{cada} elemento de $A$ tiene \textbf{exactamente una} flecha hacia $B$.}
\defi{Imagen directa}{de un grupo $E\subseteq A$: a dónde llegan las flechas que salen de $E$.}
\[f(E)=\{b\in B\mid \exists a\in E,\ b=f(a)\}.\]
\defi{Imagen inversa}{de un grupo $F\subseteq B$: los elementos de $A$ cuya flecha cae dentro de $F$.}
\[f^{-1}(F)=\{a\in A\mid f(a)\in F\}.\]
\textbf{¿Por qué funciona así?} La imagen directa va ``hacia adelante'' (se siguen las flechas); la inversa va ``hacia atrás'' (se mira de dónde vienen las flechas que caen en $F$). La imagen inversa existe siempre, aunque $f$ no tenga inversa.
\end{minipage}\hfill
\begin{minipage}[t]{0.4\textwidth}
\vspace{0pt}\centering
\begin{tikzpicture}
  \draw[oval] (-0.7,-0.5) rectangle (0.7,3.0); \node at (0,3.3) {$A$};
  \draw[oval] (2.3,-0.5) rectangle (3.7,3.0); \node at (3,3.3) {$B$};
  \fill[orange!20,rounded corners] (-0.5,1.3) rectangle (0.5,2.8); \node[orange!80!black,font=\scriptsize] at (-1.05,2.05) {$E$};
  \fill[orange!20,rounded corners] (2.5,1.8) rectangle (3.5,2.8); \node[orange!80!black,font=\scriptsize] at (4.1,2.3) {$f(E)$};
  \foreach \n/\y in {a/2.4,b/1.7,c/0.9,d/0.1}{\node[v,minimum size=5mm] (\n) at (0,\y) {\scriptsize\n};}
  \foreach \n/\y in {1/2.3,2/1.3,3/0.3}{\node[v,minimum size=5mm] (r\n) at (3,\y) {\scriptsize\n};}
  \draw[cam] (a) -- (r1); \draw[cam] (b) -- (r1);
  \draw[ar] (c) -- (r2); \draw[ar] (d) -- (r3);
\end{tikzpicture}

\smallskip
\begin{tikzpicture}
  \draw[oval] (-0.7,-0.5) rectangle (0.7,3.0); \node at (0,3.3) {$A$};
  \draw[oval] (2.3,-0.5) rectangle (3.7,3.0); \node at (3,3.3) {$B$};
  \fill[green!20,rounded corners] (2.5,-0.2) rectangle (3.5,1.7); \node[green!45!black,font=\scriptsize] at (4.0,0.7) {$F$};
  \fill[green!20,rounded corners] (-0.5,-0.2) rectangle (0.5,1.2); \node[green!45!black,font=\scriptsize] at (-1.25,0.5) {$f^{-1}(F)$};
  \foreach \n/\y in {a/2.4,b/1.7,c/0.9,d/0.1}{\node[v,minimum size=5mm] (\n) at (0,\y) {\scriptsize\n};}
  \foreach \n/\y in {1/2.3,2/1.3,3/0.3}{\node[v,minimum size=5mm] (r\n) at (3,\y) {\scriptsize\n};}
  \draw[ar] (a) -- (r1); \draw[ar] (b) -- (r1);
  \draw[gs,very thick] (c) -- (r2); \draw[gs,very thick] (d) -- (r3);
\end{tikzpicture}
\end{minipage}

\medskip
\begin{minipage}[t]{0.55\textwidth}
\defi{Inyectiva}{elementos distintos de $A$ van a lugares distintos: nunca llegan dos flechas al mismo punto.}
\[(\forall a,b\in A)\ f(a)=f(b)\Rightarrow a=b.\]
Para mostrar que \textbf{no} es inyectiva basta dar dos elementos \textbf{distintos} con la misma imagen.
\defi{Sobreyectiva}{todo elemento de $B$ recibe al menos una flecha.}
\[f(A)=B.\]
Para mostrar que \textbf{no} es sobreyectiva basta dar un $y\in B$ al que no llega ninguna flecha.\\
\defi{Biyectiva}{inyectiva y sobreyectiva; solo entonces existe la función inversa $f^{-1}$.}
\end{minipage}\hfill
\begin{minipage}[t]{0.4\textwidth}
\vspace{0pt}\centering
\begin{tikzpicture}
  \foreach \n/\y in {1/1.4,2/0.7,3/0}{\node[v,minimum size=5mm] (L\n) at (0,\y) {\scriptsize\n};}
  \foreach \n/\y in {1/1.4,2/0.7,3/0}{\node[v,minimum size=5mm] (D\n) at (1.6,\y) {\scriptsize\n};}
  \draw[bad] (L1) -- (D1); \draw[bad] (L2) -- (D1); \draw[ar] (L3) -- (D2);
  \node[red!75!black,font=\scriptsize,align=center] at (0.8,-0.7) {no inyectiva y\\ no sobreyectiva ($3$ sin flecha)};
  \begin{scope}[xshift=2.6cm]
  \foreach \n/\y in {1/1.4,2/0.7,3/0}{\node[v,minimum size=5mm] (L\n) at (0,\y) {\scriptsize\n};}
  \foreach \n/\y in {1/1.4,2/0.7,3/0}{\node[v,minimum size=5mm] (D\n) at (1.6,\y) {\scriptsize\n};}
  \draw[ok] (L1) -- (D2); \draw[ok] (L2) -- (D3); \draw[ok] (L3) -- (D1);
  \node[blue!60!black,font=\scriptsize,align=center] at (0.8,-0.7) {inyectiva y\\ sobreyectiva};
  \end{scope}
\end{tikzpicture}
\end{minipage}

% ============================================================
\newcommand{\sagitalF}[1]{%
\begin{tikzpicture}[yscale=0.62]
  \draw[oval] (-0.8,-0.7) rectangle (0.8,8.7); \node at (0,9.15) {$A\times B$};
  \draw[oval] (4.2,0.6) rectangle (5.4,7.4); \node at (4.8,7.85) {$C$};
  \foreach \a/\b/\y in {1/1/8,1/3/7,1/5/6,2/1/5,2/3/4,2/5/3,4/1/2,4/3/1,4/5/0}{%
    \node[font=\small,inner sep=1pt] (p\a\b) at (0,\y) {$(\a,\b)$};}
  \foreach \c/\y in {1/6.6,3/5.3,5/4,7/2.7,9/1.4}{\node[v,minimum size=5.5mm] (c\c) at (4.8,\y) {\c};}
  \foreach \p/\c in {11/1,13/3,15/5,21/3,23/5,25/7,41/5,43/7,45/9}{\draw[gr] (p\p.east) -- (c\c);}
  #1
\end{tikzpicture}}

\clearpage
\section*{Pregunta 5a}
\addcontentsline{toc}{section}{Pregunta 5}

\textbf{\# 5.} Considere los conjuntos $A=\{1,2,4\}$, $B=\{1,3,5\}$, y $C=\{1,3,5,7,9\}$. Se define la función $f:A\times B\to C$ por:
\[f((x,y))=\begin{cases}x+y, & \text{si } x \text{ es par},\\ xy, & \text{si } x \text{ es impar}.\end{cases}\]
a) \textbf{[2 puntos]} Calcule \hl{$f(\{1,2\}\times\{1,3\})$}.\\
b) \textbf{[2 puntos]} Calcule $f^{-1}(\{3,5\})$.\\
c) \textbf{[1 punto]} Determine si $f$ es inyectiva.\\
d) \textbf{[1 punto]} Determine si $f$ es sobreyectiva.

\begin{minipage}[c]{0.45\textwidth}
\centering
\renewcommand{\arraystretch}{1.2}
\begin{tabular}{|c|c|c|}\hline
$(x,y)$ & regla & $f((x,y))$\\\hline
$(1,1)$ & $1\cdot1$ & $1$\\
$(1,3)$ & $1\cdot3$ & $3$\\
$(1,5)$ & $1\cdot5$ & $5$\\
$(2,1)$ & $2+1$ & $3$\\
$(2,3)$ & $2+3$ & $5$\\
$(2,5)$ & $2+5$ & $7$\\
$(4,1)$ & $4+1$ & $5$\\
$(4,3)$ & $4+3$ & $7$\\
$(4,5)$ & $4+5$ & $9$\\\hline
\end{tabular}
\end{minipage}\hfill
\begin{minipage}[c]{0.5\textwidth}
\centering
\sagitalF{\foreach \p/\c in {11/1,13/3,21/3,23/5}{\draw[cam] (p\p.east) -- (c\c);}
  \fill[orange!25,opacity=0.6,rounded corners] (-0.65,3.6) rectangle (0.65,5.4);
  \fill[orange!25,opacity=0.6,rounded corners] (-0.65,6.6) rectangle (0.65,8.4);}
\end{minipage}

\calculo{f(\{1,2\}\times\{1,3\})=\{\bajo{1}{f(1,1)},\ \bajo{3}{f(1,3)=f(2,1)},\ \bajo{5}{f(2,3)}\}}

\begin{explica}
\textbf{Paso 1, tabla de valores:} el dominio es $A\times B$, que tiene $3\cdot3=9$ pares. Para cada par se mira si $x$ es par o impar: $x=1$ es impar, así que se \textbf{multiplica} $x\cdot y$; $x=2$ y $x=4$ son pares, así que se \textbf{suma} $x+y$. Así sale la tabla de la izquierda y el diagrama sagital.

\textbf{Paso 2, el conjunto de salida:} por Def. de producto cartesiano, $\{1,2\}\times\{1,3\}=\{(1,1),(1,3),(2,1),(2,3)\}$ (zonas naranjas).

\textbf{Paso 3, imagen directa:} por Def. de imagen directa (Teoría 5) se siguen las flechas que salen de esos cuatro pares: $f(1,1)=1$, $f(1,3)=3$, $f(2,1)=3$, $f(2,3)=5$.

\textbf{Paso 4, se junta sin repetir:} el $3$ aparece dos veces pero en un conjunto se escribe una sola vez: $\{1,3,5\}$.

\textbf{¿Tiene sentido?} Son a lo sumo $4$ valores y salieron $3$ porque dos pares caen en el mismo $3$.
\end{explica}

\textbf{Respuesta:} $f(\{1,2\}\times\{1,3\})=\{1,3,5\}$.

% ============================================================
\clearpage
\section*{Pregunta 5b}

b) \textbf{[2 puntos]} Calcule \hl{$f^{-1}(\{3,5\})$}.

\begin{center}
\sagitalF{\foreach \p/\c in {13/3,21/3,15/5,23/5,41/5}{\draw[gs,very thick] (p\p.east) -- (c\c);}
  \draw[green!45!black,very thick,rounded corners] (4.4,3.5) rectangle (5.2,5.8);
  \node[green!45!black,font=\small] at (6.2,4.65) {$F=\{3,5\}$};}
\end{center}

\calculo{f^{-1}(\{3,5\})=\{\bajo{(1,3)}{\to3},\ \bajo{(1,5)}{\to5},\ \bajo{(2,1)}{\to3},\ \bajo{(2,3)}{\to5},\ \bajo{(4,1)}{\to5}\}}

\begin{explica}
\textbf{Definición que se usa:} imagen inversa (Teoría 5): $f^{-1}(F)=\{(x,y)\in A\times B\mid f((x,y))\in F\}$. Se buscan los pares cuya flecha \textbf{cae} en $3$ o en $5$.

\textbf{Paso 1, los que caen en $3$:} en la tabla de la Pregunta 5a, la columna $f$ vale $3$ en $(1,3)$ y $(2,1)$.

\textbf{Paso 2, los que caen en $5$:} vale $5$ en $(1,5)$, $(2,3)$ y $(4,1)$.

\textbf{Paso 3, se juntan:} $2+3=5$ pares en total (flechas verdes).

\textbf{¿Tiene sentido?} Los otros cuatro pares, $(1,1)\to1$, $(2,5)\to7$, $(4,3)\to7$, $(4,5)\to9$, caen fuera de $\{3,5\}$, así que no entran. Las $9$ filas de la tabla quedan revisadas: $5+4=9$.
\end{explica}

\textbf{Respuesta:} $f^{-1}(\{3,5\})=\{(1,3),(1,5),(2,1),(2,3),(4,1)\}$.

% ============================================================
\clearpage
\section*{Pregunta 5c y 5d}

c) \textbf{[1 punto]} Determine si $f$ es \hl{inyectiva}.\\
d) \textbf{[1 punto]} Determine si $f$ es \hl{sobreyectiva}.

\begin{center}
\sagitalF{\draw[bad] (p13.east) -- (c3); \draw[bad] (p21.east) -- (c3);
  \foreach \c in {1,3,5,7,9}{\draw[green!45!black,thick] (c\c) circle (3.6mm);}
  \node[red!75!black,font=\scriptsize,align=left,anchor=west] at (5.6,5.3) {dos flechas\\ llegan a $3$};
  \node[green!45!black,font=\scriptsize,align=left,anchor=west] at (5.6,2.0) {todos los de $C$\\ reciben flecha};}
\end{center}

\calculo{\bajo{f((1,3))}{1\cdot3}=\bajo{f((2,1))}{2+1}=3,\qquad f[A\times B]=\{1,3,5,7,9\}=C}

\begin{explica}
\textbf{c) Inyectiva, definición:} (Teoría 5) $f(a)=f(b)\Rightarrow a=b$: no pueden llegar dos flechas al mismo punto. Para negarlo basta un ejemplo.

\textbf{Paso 1:} en la tabla, el valor $3$ aparece dos veces: $f((1,3))=1\cdot3=3$ (porque $1$ es impar) y $f((2,1))=2+1=3$ (porque $2$ es par). Los pares $(1,3)$ y $(2,1)$ son distintos, pero tienen la misma imagen: no es inyectiva.

\textbf{d) Sobreyectiva, definición:} $f(A\times B)=C$: cada elemento de $C$ debe recibir al menos una flecha.

\textbf{Paso 2:} se revisa cada elemento de $C$ con la tabla: $1=f((1,1))$, $3=f((1,3))$, $5=f((1,5))$, $7=f((2,5))$, $9=f((4,5))$. Los cinco reciben flecha, así que el conjunto de imágenes es todo $C$.

\textbf{¿Tiene sentido?} El dominio tiene $9$ elementos y $C$ solo $5$: con $9$ flechas hacia $5$ puntos es imposible que no se repita ningún punto, por eso la inyectividad tenía que fallar; en cambio, sí alcanzan para cubrir los $5$.
\end{explica}

c) La función no es inyectiva, pues $f((1,3))=f((2,1))=3$, con elementos distintos del dominio.\\
d) La función sí es sobreyectiva, ya que $f[A\times B]=\{1,3,5,7,9\}=C$.

\textbf{Respuesta:} c) $f$ no es inyectiva. d) $f$ sí es sobreyectiva.

% ============================================================
\clearpage
\section*{Pregunta 6}
\addcontentsline{toc}{section}{Pregunta 6}

\textbf{\# 6. [3 puntos]} Sean $f:A\to B$ una función y $E\subseteq A$ y $F\subseteq B$. Demuestre que:
\[\colorbox{green!35}{$E\cap f^{-1}(F)\subseteq f^{-1}(f(E)\cap F)$}.\]

\begin{center}
\begin{tikzpicture}
  \draw[oval] (-1.6,-1.5) rectangle (1.6,1.6); \node at (0,1.95) {$A$};
  \draw[oval] (4.4,-1.5) rectangle (7.6,1.6); \node at (6,1.95) {$B$};
  \draw[blue!60!black,thick,fill=blue!8,fill opacity=0.6] (-0.45,0.1) circle (0.95); \node[blue!60!black,font=\small] at (-1.2,1.05) {$E$};
  \draw[green!45!black,thick,fill=green!10,fill opacity=0.6] (0.45,-0.2) circle (0.95); \node[green!45!black,font=\scriptsize] at (0.75,-1.32) {$f^{-1}(F)$};
  \draw[blue!60!black,thick,fill=blue!8,fill opacity=0.6] (5.55,0.1) circle (0.95); \node[blue!60!black,font=\small] at (4.75,1.05) {$f(E)$};
  \draw[green!45!black,thick,fill=green!10,fill opacity=0.6] (6.45,-0.2) circle (0.95); \node[green!45!black,font=\small] at (7.25,-1.15) {$F$};
  \fill (0,-0.05) circle (2pt) node[left,font=\small] {$x$};
  \fill (6,-0.05) circle (2pt) node[right,font=\small] {$f(x)$};
  \draw[cam] (0.1,-0.05) to[bend left=20] node[above,font=\small] {$f$} (5.9,-0.05);
\end{tikzpicture}
\end{center}

\calculo{x\in\bajo{E}{\Rightarrow f(x)\in f(E)}\ \wedge\ x\in\bajo{f^{-1}(F)}{\Rightarrow f(x)\in F}\ \Longrightarrow\ f(x)\in f(E)\cap F}

\begin{explica}
\textbf{Qué significa $\subseteq$:} todo elemento del conjunto de la izquierda también está en el de la derecha. Se toma un $x$ fijo pero arbitrario de la izquierda y se muestra que está a la derecha.

\textbf{Paso 1, Def. $\cap$:} $x\in E\cap f^{-1}(F)$ quiere decir que $x\in E$ \textbf{y} $x\in f^{-1}(F)$.

\textbf{Paso 2, Def. imagen inversa:} $x\in f^{-1}(F)$ quiere decir que la flecha de $x$ cae en $F$: $f(x)\in F$.

\textbf{Paso 3, Def. imagen directa:} como $x\in E$, su imagen $f(x)$ es una de las que forman $f(E)$: $f(x)\in f(E)$.

\textbf{Paso 4, Def. $\cap$ otra vez:} $f(x)$ está en $f(E)$ y en $F$, o sea en $f(E)\cap F$ (zona donde se cruzan los círculos de $B$).

\textbf{Paso 5, Def. imagen inversa otra vez:} si la flecha de $x$ cae en $f(E)\cap F$, entonces $x\in f^{-1}(f(E)\cap F)$.

\textbf{¿Tiene sentido?} En el dibujo, $x$ está en el cruce de los círculos de $A$ y su flecha llega al cruce de los círculos de $B$.
\end{explica}

\textbf{HQD:} $(\forall x\in A)\big[x\in E\cap f^{-1}(F)\Rightarrow x\in f^{-1}(f(E)\cap F)\big]$.

Sea $x\in A$ fijo pero arbitrario tal que $x\in E\cap f^{-1}(F)$. Entonces:
\begin{align*}
x\in E\cap f^{-1}(F)&\Rightarrow x\in E\wedge x\in f^{-1}(F)\jus{Def. $\cap$}\\
&\Rightarrow x\in E\wedge f(x)\in F\jus{Def. imagen inversa}\\
&\Rightarrow f(x)\in f(E)\wedge f(x)\in F\jus{Def. imagen directa}\\
&\Rightarrow f(x)\in f(E)\cap F\jus{Def. $\cap$}\\
&\Rightarrow x\in f^{-1}(f(E)\cap F)\jus{Def. imagen inversa}
\end{align*}
Por lo tanto, $E\cap f^{-1}(F)\subseteq f^{-1}(f(E)\cap F)$.

\textbf{Respuesta:} queda demostrada la inclusión.

% ============================================================
\clearpage
\section*{Pregunta 7a}
\addcontentsline{toc}{section}{Pregunta 7}

\textbf{\# 7.} Sea $f:\mathbb R-\{-2\}\to\mathbb R$ la función definida por:
\[f(x)=\frac{3x}{x+2}+1.\]
a) Demuestre que $f$ es \hl{inyectiva}. \hfill\textbf{[2 puntos]}\\
b) Justifique que $f$ \hl{no es sobreyectiva}. \hfill\textbf{[2 puntos]}

\begin{center}
\begin{tikzpicture}
  \node[v,minimum size=8mm] (a) at (0,0.8) {$x_1$}; \node[v,minimum size=8mm] (b) at (0,-0.8) {$x_2$};
  \node[v,minimum size=8mm] (y) at (3.5,0) {$y$};
  \draw[ar] (a) -- (y); \draw[ar] (b) -- (y);
  \node[font=\small,align=left,anchor=west] at (4.3,0) {si llegan al mismo $y$,\\ la cuenta obliga a $x_1=x_2$:\\ en realidad es un solo punto};
\end{tikzpicture}
\end{center}

\calculo{\frac{3x_1}{x_1+2}=\frac{3x_2}{x_2+2}\ \Rightarrow\ \bajo{x_1x_2+2x_1=x_1x_2+2x_2}{\text{se cancela }x_1x_2}\ \Rightarrow\ x_1=x_2}

\begin{explica}
\textbf{Definición que se usa:} inyectiva (Teoría 5): $f(x_1)=f(x_2)\Rightarrow x_1=x_2$. Se suponen dos entradas con la misma salida y, con álgebra, se llega a que son la misma.

\textbf{Paso 1, Def. de $f$:} se escribe $f(x_1)$ y $f(x_2)$ con la fórmula y se igualan.

\textbf{Paso 2, se quita el $+1$:} se resta $1$ a ambos lados.

\textbf{Paso 3, se quitan las fracciones:} se multiplica en cruz. Se puede porque $x_1+2\neq0$ y $x_2+2\neq0$ (el $-2$ no está en el dominio). Queda $3x_1(x_2+2)=3x_2(x_1+2)$ y se divide entre $3$.

\textbf{Paso 4, Dist.:} $x_1(x_2+2)=x_1x_2+2x_1$ y $x_2(x_1+2)=x_1x_2+2x_2$.

\textbf{Paso 5, Cancelación:} $x_1x_2$ está a ambos lados y se cancela; queda $2x_1=2x_2$ y, dividiendo entre $2$, $x_1=x_2$.

\textbf{¿Tiene sentido?} Por ejemplo $f(0)=1$ y $f(1)=2$: entradas distintas dieron salidas distintas, como se espera.
\end{explica}

\textbf{HQD:} $(\forall x_1,x_2\in\mathbb R-\{-2\})\big[f(x_1)=f(x_2)\Rightarrow x_1=x_2\big]$.

Sean $x_1,x_2\in\mathbb R-\{-2\}$ fijos pero arbitrarios tales que $f(x_1)=f(x_2)$. Entonces:
\begin{align*}
f(x_1)=f(x_2)&\Rightarrow\frac{3x_1}{x_1+2}+1=\frac{3x_2}{x_2+2}+1\jus{Def. de $f$}\\
&\Rightarrow\frac{3x_1}{x_1+2}=\frac{3x_2}{x_2+2}\jus{Resta de $1$}\\
&\Rightarrow x_1(x_2+2)=x_2(x_1+2)\jus{$x_1+2\neq0,\ x_2+2\neq0$}\\
&\Rightarrow x_1x_2+2x_1=x_1x_2+2x_2\jus{Dist.}\\
&\Rightarrow 2x_1=2x_2\jus{Cancelación}\\
&\Rightarrow x_1=x_2.
\end{align*}
Por lo tanto, $f$ es inyectiva.

\textbf{Respuesta:} $f$ es inyectiva.

% ============================================================
\clearpage
\section*{Pregunta 7b}

b) Justifique que $f$ \hl{no es sobreyectiva}. \hfill\textbf{[2 puntos]}

\begin{center}
\begin{tikzpicture}
\begin{axis}[width=0.62\textwidth,height=5.2cm,axis lines=middle,xmin=-9,xmax=7,ymin=-5,ymax=13,
  xtick={-2},ytick={4},xlabel={$x$},ylabel={$y$},samples=120,clip=true,
  every axis x label/.style={at={(ticklabel* cs:1)},anchor=west},
  every axis y label/.style={at={(ticklabel* cs:1)},anchor=south}]
  \addplot[blue!60!black,very thick,domain=-9:-2.6] {3*x/(x+2)+1};
  \addplot[blue!60!black,very thick,domain=-1.55:7] {3*x/(x+2)+1};
  \addplot[red!75!black,very thick,dashed,domain=-9:7] {4};
  \draw[gray,dashed] (axis cs:-2,-5) -- (axis cs:-2,13);
  \node[red!75!black,font=\scriptsize\bfseries,anchor=south west] at (axis cs:1.5,4.3) {$y=4$: nunca se alcanza};
\end{axis}
\end{tikzpicture}
\end{center}

\calculo{x=\frac{2-2y}{\bajo{y-4}{=0\ \text{si}\ y=4}}\qquad f(x)=4\Rightarrow 3x=3x+6\Rightarrow \bajo{0=6}{\text{absurdo}}}

\begin{explica}
\textbf{Definición que se usa:} sobreyectiva (Teoría 5): todo $y\in\mathbb R$ debe tener algún $x\in\mathbb R-\{-2\}$ con $f(x)=y$. Para negarlo basta un $y$ sin preimagen.

\textbf{Paso 1, despejar $x$:} se iguala $f(x)=y$. Se resta $1$: $y-1=\frac{3x}{x+2}$. Se multiplica por $x+2$: $(y-1)(x+2)=3x$. Se distribuye: $xy-x+2y-2=3x$. Se pasan los términos con $x$ a la izquierda y los demás a la derecha: $xy-4x=2-2y$. Se factoriza la $x$: $x(y-4)=2-2y$. Se divide: $x=\frac{2-2y}{y-4}$.

\textbf{Paso 2, ¿dónde falla?} No se puede dividir entre cero, así que la fórmula no da ningún $x$ cuando $y-4=0$, o sea $y=4$.

\textbf{Paso 3, comprobación directa:} si existiera $x$ con $f(x)=4$: $\frac{3x}{x+2}+1=4\Rightarrow\frac{3x}{x+2}=3\Rightarrow 3x=3(x+2)=3x+6\Rightarrow 0=6$, que es falso. Entonces ningún $x$ del dominio da $4$.

\textbf{¿Tiene sentido?} En la gráfica la curva se acerca a la línea roja $y=4$ pero nunca la toca. El $4$ es un número real (está en el codominio $\mathbb R$) y no recibe flecha: $f(\mathbb R-\{-2\})\neq\mathbb R$.
\end{explica}

Supóngase que $f$ fuera sobreyectiva. Entonces, para todo $y\in\mathbb R$, debería existir $x\in\mathbb R-\{-2\}$ tal que $f(x)=y$. Se despeja $x$ en términos de $y$:
\begin{align*}
y&=\frac{3x}{x+2}+1\\
y-1&=\frac{3x}{x+2}\\
(y-1)(x+2)&=3x\\
xy-x+2y-2&=3x\\
xy-4x&=2-2y\\
x(y-4)&=2-2y\\
x&=\frac{2-2y}{y-4}.
\end{align*}
Esta expresión no está definida cuando $y-4=0$, es decir, cuando $y=4$. En efecto, $f(x)=4\Rightarrow\frac{3x}{x+2}=3\Rightarrow3x=3x+6\Rightarrow0=6$, absurdo. Por lo tanto, para $y=4$ no existe $x\in\mathbb R-\{-2\}$ tal que $f(x)=4$. Como $4\in\mathbb R$, se concluye que $f$ no es sobreyectiva.

\textbf{Respuesta:} $f$ no es sobreyectiva: $4\in\mathbb R$ no tiene preimagen.

\end{document}
